\documentclass{beamer}


%%%%%%%%%%%%%%%%%%%%%%%%%%%%%%%%%%%%%%%%%%%%%%%%%%%%%%%%%%%%%
%% Theme Packages %%%%%%%%%%%%%%%%%%%%%%%%%%%%%%%%%%%%%%%%%%%
%%%%%%%%%%%%%%%%%%%%%%%%%%%%%%%%%%%%%%%%%%%%%%%%%%%%%%%%%%%%%
\usepackage{environ}
\usepackage[many]{tcolorbox}
\usepackage{tikz}
\usetikzlibrary{calc,positioning,shadows,shapes}
\usepackage[absolute,overlay]{textpos}
\usepackage{xifthen}
\usepackage{adjustbox}
\usepackage{varwidth}
\usepackage{transparent}

%%%%%%%%%%%%%%%%%%%%%%%%%%%%%%%%%%%%%%%%%%%%%%%%%%%%%%%%%%%%%
%% Recommended Packages %%%%%%%%%%%%%%%%%%%%%%%%%%%%%%%%%%%%%
%%%%%%%%%%%%%%%%%%%%%%%%%%%%%%%%%%%%%%%%%%%%%%%%%%%%%%%%%%%%%
% \usepackage{natbib} 
% \usepackage[english]{babel} % english language
% \usepackage[utf8]{inputenc} % utf8 encoding of input files
% \usepackage[T1]{fontenc} % T1 encoding of output pdf files
% \usepackage{mathtools} % mathematical notation
% \usepackage{mathpartir} % automatical alignment of multiple entries
% \usepackage{multicol} % multiple columns layout
% \usepackage{graphicx} % better handling of images
% \usepackage{relsize} % for \smaller, ...
% \usepackage{fontawesome} % for icons
% \usepackage{bussproofs} % for proof trees
% \usepackage{stackengine} % for stacking objects
% \usepackage{changepage} % adjust page layout
%%%%%%%%%%%%%%%%%%%%%%%%%%%%%%%%%%%%%%%%%%%%%%%%%%%%%%%%%%%%%
%% Variables %%%%%%%%%%%%%%%%%%%%%%%%%%%%%%%%%%%%%%%%%%%%%%%%
%%%%%%%%%%%%%%%%%%%%%%%%%%%%%%%%%%%%%%%%%%%%%%%%%%%%%%%%%%%%%

\newcommand{\primarycolor}[1]{\def\theprimarycolor{#1}}
\newcommand{\secondarycolor}[1]{\def\thesecondarycolor{#1}}
\newcommand{\counterbackground}[1]{\def\thecounterbackground{#1}}

%%%%%%%%%%%%%%%%%%%%%%%%%%%%%%%%%%%%%%%%%%%%%%%%%%%%%%%%%%%%%
%% Colors %%%%%%%%%%%%%%%%%%%%%%%%%%%%%%%%%%%%%%%%%%%%%%%%%%%
%%%%%%%%%%%%%%%%%%%%%%%%%%%%%%%%%%%%%%%%%%%%%%%%%%%%%%%%%%%%%

\definecolor{darkcerulean}{rgb}{0.063661, 0.257392, 0.477463}
\definecolor{scooter}{rgb}{0.161162, 0.775760, 0.885416}
\definecolor{fireenginered}{rgb}{0.757398, 0.151664, 0.177088}
\definecolor{pomegranate}{rgb}{0.927557, 0.260794, 0.138377}
\definecolor{mineshaft}{rgb}{0.199974, 0.200015, 0.199971}
\definecolor{tingray}{rgb}{0.501904, 0.501994, 0.501898}


\primarycolor{darkcerulean}
\secondarycolor{scooter}

%%%%%%%%%%%%%%%%%%%%%%%%%%%%%%%%%%%%%%%%%%%%%%%%%%%%%%%%%%%%%
%% Beamer Layout %%%%%%%%%%%%%%%%%%%%%%%%%%%%%%%%%%%%%%%%%%%%
%%%%%%%%%%%%%%%%%%%%%%%%%%%%%%%%%%%%%%%%%%%%%%%%%%%%%%%%%%%%%
\useoutertheme[subsection=false]{miniframes}
\useinnertheme{default}
\usefonttheme[onlymath]{serif}
\setbeamercolor*{normal text}{fg=black,bg=white}
\setbeamercolor*{section in head/foot}{fg=white, bg=\theprimarycolor}
\setbeamercolor*{frametitle}{fg=\theprimarycolor, bg=white}
\setbeamercolor*{titlelike}{fg=\theprimarycolor,bg=white}
\setbeamercolor*{structure}{fg=\theprimarycolor}
 
\setbeamertemplate{logo}{
    \vspace{-1em}
    \begin{tikzpicture}
        \node {\thecounterbackground};
        \node {{\tiny\color{white}\insertframenumber}};
    \end{tikzpicture}
}
\setbeamertemplate{footline}{}
\beamertemplatenavigationsymbolsempty

\setbeamertemplate{itemize items}{
    \begin{tikzpicture}
        \fill (0,0) circle (2pt);
        \draw (0,0) circle (3pt);
    \end{tikzpicture}
}

\setbeamertemplate{enumerate items}[default]


%****************************************************************************
% The following is a hack to allow the transparent package to be loaded at the
% same time as the tikz package: 
% https://tex.stackexchange.com/questions/253401/tikz-and-ctable-incompatibility-gives-error-when-printing/253417#253417


\makeatletter
\begingroup\expandafter\expandafter\expandafter\endgroup
\expandafter\ifx\csname pgfutil@addpdfresource@extgs\endcsname\relax
\else
  \AtBeginDocument{%
    % \pgf@sys@addpdfresource@extgs@plain{%
    \pgfutil@addpdfresource@extgs{%
      \TRP@list
    }%
  }%  
  \let\TRP@addresource\relax
\fi
\makeatother
%****************************************************************************




%***************************************************************************************************
% The following magic makes sure that the pnote command produces notes that are displayed in pdfpc:

% create a new file handle
\newwrite\pdfpcnotesfile

% open file on \begin{document}
\AtBeginDocument{%
	\immediate\openout\pdfpcnotesfile\jobname.pdfpc\relax
	\immediate\write\pdfpcnotesfile{[notes]}
}
% define a # http://tex.stackexchange.com/a/37757/10327
\begingroup
	\catcode`\#=12
	\gdef\hashchar{#}%
\endgroup


\def\lastframenumber{0}

% define command \pnote{} that works like note but
% additionally writes notes to file in pdfpc readable format
\newcommand{\pnote}[1]{%
	% keep normal notes working
	\note{#1}%

	% if frame changed - write a new header
	\ifdim\theframenumber pt>\lastframenumber pt
		\let\lastframenumber\theframenumber
		\begingroup
			\let\#\hashchar
			\immediate\write\pdfpcnotesfile{\#\#\# \theframenumber}%
		\endgroup
	\fi

	% write note to file
	\immediate\write\pdfpcnotesfile{\unexpanded{#1}}%
}
% close file on \begin{document}
\AtEndDocument{%
	\immediate\closeout\pdfpcnotesfile
}

%            END OF pdfpc MAGIC
%***************************************************************************************************

%%%%%%%%%%%%%%%%%%%%%%%%%%%%%%%%%%%%%%%%%%%%%%%%%%%%%%%%%%%%%
%%% Boxes
%%%%%%%%%%%%%%%%%%%%%%%%%%%%%%%%%%%%%%%%%%%%%%%%%%%%%%%%%%%%%


% #1 – box name
% #2 – left color
% #3 – right color
% #4 – Title (may contain #1)
\newcommand{\newlargebox}[4]{
    \newtcolorbox{#1}[1][]{
        beamer,
        width=\textwidth,
        titlerule=3mm,
        colframe=white,
        coltitle={#2},
        bottom=4pt, 
        top=-4pt,
        left=6pt,   
        right=6pt,   
        toptitle=4pt,
        bottomtitle=-2pt,
        fonttitle=\bfseries\large,
        adjusted title={#4},
        outer arc=.5mm,
        arc=.5mm,
        no shadow,
        fuzzy shadow={1mm}{-1mm}{-1.2mm}{.7mm}{black!20},
        interior titled code={
          \path [left color = {#2}, right color = {#3}]
          (title.south west) + (8pt, 0) rectangle ++(\textwidth-8pt, 0.02);
        }
      }
}


\newtcolorbox{emptybox}[1][]{
    beamer,
    width=\textwidth,
    % enlarge left by=-3pt,
    titlerule=0mm,
    colframe=white,
    coltitle=black,
    bottom=8pt, 
    top=-10pt,
    left=6pt,  
    notitle, 
    adjusted title={},
    outer arc=.5mm,
    arc=.5mm,
    no shadow,
    fuzzy shadow={1mm}{-1mm}{-1.2mm}{.7mm}{black!20},
    interior titled code={}
}

\newtcolorbox{plainbox}[1][]{
    beamer,
    width=\textwidth,
    % enlarge left by=-3pt,
    titlerule=0mm,
    colframe=white,
    coltitle=black,
    bottom=8pt,
    toptitle=4pt, 
    top=2pt,
    left=6pt,  
    notitle, 
    adjusted title={#1},
    outer arc=.5mm,
    arc=.5mm,
    no shadow,
    fuzzy shadow={1mm}{-1mm}{-1.2mm}{.7mm}{black!20},
    interior titled code={}
}

\newlargebox{defaultbox}{\theprimarycolor}{\thesecondarycolor}{#1} 
\newlargebox{darkbox}{mineshaft}{tingray}{#1} 
\newlargebox{graybox}{tingray}{white}{#1} 

\newlargebox{theorembox}{\theprimarycolor}{\thesecondarycolor}{Theorem {\ifthenelse{\equal{#1}{}}{}{(#1)}}} 
\newlargebox{definitionbox}{\theprimarycolor}{\thesecondarycolor}{Definition {\ifthenelse{\equal{#1}{}}{}{(#1)}}} 
\newlargebox{proofbox}{mineshaft}{tingray}{Proof{\ifthenelse{\equal{#1}{}}{}{ #1}}.} 

\newcommand{\smallbox}[1]{\begin{emptybox}\textbf{#1}\end{emptybox}}
\newcommand{\defaultsmallbox}[1]{\smallbox{\color{\theprimarycolor}#1}}
\newcommand{\graysmallbox}[1]{\smallbox{\color{tingray}#1}}


%%%%%%%%%%%%%%%%%%%%%%%%%%%%%%%%%%%%%%%%%%%%%%%%%%%%%%%%%%%%%
%% Common Latex Definitions %%%%%%%%%%%%%%%%%%%%%%%%%%%%%%%%%
%%%%%%%%%%%%%%%%%%%%%%%%%%%%%%%%%%%%%%%%%%%%%%%%%%%%%%%%%%%%%

\renewcommand{\iff}{\text{iff}}
\newcommand{\set}[1]{\ensuremath{\left\{#1\right\}}}
\newcommand{\bnfor}{~|~}
\newcommand{\brackets}[1]{\left[{#1}\right]}
\newcommand{\angles}[1]{\left\langle{#1}\right\rangle}
\newcommand{\round}[1]{\left({#1}\right)}
\newcommand{\cev}[1]{\reflectbox{\ensuremath{\vec{\reflectbox{\ensuremath{#1}}}}}} % reversed vec
\newcommand{\abs}[1]{\left|#1\right|}
\newcommand{\semantics}[2][]{\left\llbracket#2\ifthenelse{\isempty{#1}}{\right\rrbracket}{\right\rrbracket_{#1}}}
\newcommand{\powerset}[1]{\mathcal{P}(#1)}
\newcommand{\trans}[1]{\overset{#1}\to}

% Number Sets
\newcommand{\BB}{\ensuremath{\mathbb{B}}}
\newcommand{\NN}{\ensuremath{\mathbb{N}}}
\newcommand{\ZZ}{\ensuremath{\mathbb{Z}}}
\newcommand{\QQ}{\ensuremath{\mathbb{Q}}}
\newcommand{\RR}{\ensuremath{\mathbb{R}}}
\newcommand{\CC}{\ensuremath{\mathbb{C}}}


\title{Beamer Theme Example}
\author{Simon Spies}

\begin{document}

\begin{frame}[plain]
    \maketitle
\end{frame} 

\section{Basics}
\subsection{}

\begin{frame}{Colors}
    The theme mainly relies on two colors, the primary color \tikz{\fill[fill=\theprimarycolor] (0,0) rectangle ++(10pt,10pt)}
    and the secondary color \tikz{\fill[fill=\thesecondarycolor] (0,0) rectangle ++(10pt,10pt)}. For example, the secondary color 
    is used to \highlight{highlight text}.\\[2em]

\end{frame}


{
    \primarycolor{red}
    \secondarycolor{green}
    \begin{frame}[fragile]{Colors}
        The theme mainly relies on two colors, the primary color \tikz{\fill[fill=\theprimarycolor] (0,0) rectangle ++(10pt,10pt)}
        and the secondary color \tikz{\fill[fill=\thesecondarycolor] (0,0) rectangle ++(10pt,10pt)}. For example, the secondary color 
        is used to \highlight{highlight text}.\\[2em]

        To change the primary color to \texttt{red} and the secondary color to \verb|green|,
        use \verb|\primarycolor{red}| and \verb|\secondarycolor{green}|. 
    \end{frame}
}

\begin{frame}{Lists}
    \textbf{Itemize}
    \begin{itemize}
        \item First Item
        \item Second Item
    \end{itemize}
    
    \vspace{2em}
    \textbf{Enumerate}
    \begin{enumerate}
        \item First Item
        \item Second Item
    \end{enumerate}
\end{frame}
    
\begin{frame}[fragile]{Highlighting}
    You can highlight \highlight{text} with \verb|\highlight{text}|. \verb|\highlight| is compatible 
    with beamer overlays. Thus, to only highlight \highlight<2>{text} on the second slide you can 
    simply write \verb|\highlight<2>{text}|.\\[2em]


    To change the highlighting color of the \highlight[red]{text}, pass the color 
    as an optional argument \verb|\highlight[red]{text}|.
\end{frame}

\begin{frame}[fragile]{Frame Numbers}
    The color of the frame numbers can be controlled with 
    \verb|\framenumbercolor{red}|. The background with
    \verb|\framenumberbackground{...}|.
\end{frame}

{
    \framenumbercolor{\thesecondarycolor}
    \framenumberbackground
    {{\color{\theprimarycolor}\huge$\blacksquare$}}
\begin{frame}[fragile]{Frame Numbers}
    The color of the frame numbers can be controlled with 
    \verb|\framenumbercolor|. The background with
    \verb|\framenumberbackground|.\vspace{2em}

    For example, we can replace the background with a heart and 
    make the counter white:
    \begin{verbatim}
\framenumbercolor{\thesecondarycolor}
\framenumberbackground
    {{\color{\theprimarycolor}\huge$\blacksquare$}}
    \end{verbatim}
\end{frame}
}

\section{Boxes}
\subsection{}

\begin{frame}{Color Boxes}
    \begin{defaultbox}[Default Box]
        This is a simple box.
    \end{defaultbox}

    \begin{darkbox}[Dark Box]
        This is a dark box.
    \end{darkbox}

    \begin{graybox}[Gray Box]
        This is a gray box.
    \end{graybox}
\end{frame}

\begin{frame}{Other Boxes}
    \begin{theorembox}[Really True]
        \[\mathsf{true} = \mathsf{false}\]
    \end{theorembox}

    \begin{definitionbox}[Something Cool]
        This is a great definition.
    \end{definitionbox}

    \begin{proofbox}
        This is really a proof.
    \end{proofbox}
\end{frame}

\begin{frame}{Small Boxes}
    \defaultsmallbox{Default}
    \smallbox{Small}
    \graysmallbox{Gray}
\end{frame}

\begin{frame}{Simple Boxes}
    \begin{plainbox}[Plain Box]
        This is a plain box.
    \end{plainbox}

    \begin{emptybox}
        This is an empty box.
    \end{emptybox}
\end{frame}

\section{Timeline}
\subsection{} 




\begin{frame}{Timeline}
    
    \begin{textblock*}{\paperwidth} (0cm,5pt)
        \begin{center}
            \begin{tikzpicture}
                \coordinate (top) at (0,0);
                \coordinate[below=8 of top] (bottom);
                \coordinate[below=2 of top] (one);
                \coordinate[below=3.5 of top] (two);
                \coordinate[below=5 of top] (three);
                \coordinate[below=6.5 of top] (four);
    
                \draw[->, \theprimarycolor] (top) -- (bottom);
    
                \visible<2->{
                    \begin{timelinenode}[right]{one}{.4\linewidth}{First Item}
                        Text here.
                    \end{timelinenode}
                }
                \visible<3->{
                    \begin{timelinenode}[left]{two}{.4\linewidth}{Second Item}
                        Text here.
                    \end{timelinenode}
                }
                \visible<4->{
                    \begin{timelinenode}[right]{three}{.4\linewidth}{Third Item}
                        Text here.
                    \end{timelinenode}
                }
                \visible<5->{
                    \begin{graytimelinenode}[left]{four}{.4\linewidth}{Fourth Item}
                        Text here.
                    \end{graytimelinenode}
                }
            \end{tikzpicture}        
        \end{center}
    \end{textblock*}
\end{frame}


\begin{frame}{Simple Timeline}
    \begin{textblock*}{\paperwidth} (0cm,5pt)
        \begin{center}
            \begin{tikzpicture}
                \coordinate (top) at (0,0);
                \coordinate[below=8 of top] (bottom);
                \coordinate[below=2 of top] (one);
                \coordinate[below=3.5 of top] (two);
                \coordinate[below=5 of top] (three);
                \coordinate[below=6.5 of top] (four);
    
                \draw[->, \theprimarycolor] (top) -- (bottom);
    
                \smalltimelinenode[right]{one}{.2\textwidth}{First Item};
                \smalltimelinenode[left]{two}{.3\textwidth}{Second Item};
                \smalltimelinenode[right]{three}{.2\textwidth}{Third Item};
                \graysmalltimelinenode[left]{four}{.3\textwidth}{Fourth Item};
            \end{tikzpicture}        
        \end{center}
    \end{textblock*}
\end{frame}

\end{document}